% Part: computability
% Chapter: computability-theory
% Section: total

\documentclass[../../../include/open-logic-section]{subfiles}

\begin{document}

\olfileid{cmp}{thy}{tot}
\olsection{সর্বত্র-সংজ্ঞায়িততা অনির্ণেয়}

কোনো কিছু অগণনসাধ্য দেখাতে $s$-$m$-$n$ উপপাদ্যের ব্যবহারের আরও একটি
উদাহরণ বিবেচনা করা যাক। সর্বত্র-সংজ্ঞায়িত গণনসাধ্য অপেক্ষকগুলির
সূচকসমূহের সেটকে~$\fn{Tot}$ ধরি, অর্থাৎ
\[
\fn{Tot} = \Setabs{x}{\text{প্রত্যেক $y$-এর জন্য $\cfind{x}(y)\fdefined$}}.
\]

\begin{prop}
\ollabel{prop:total}
$\fn{Tot}$ গণনসাধ্য নয়।
\end{prop}

\begin{proof}
$\fn{Tot}$ গণনসাধ্য নয় দেখতে $K$ যে তার মধ্যে হ্রাসযোগ্য তা দেখালেই
যথেষ্ট। $h(x,y)$-কে সংজ্ঞায়িত করি:
\[
h(x,y) \simeq
\begin{cases}
0 & \text{যদি $x \in K$} \\
\fundefined & \text{অন্যথায়}
\end{cases}
\]
লক্ষ করো, $h(x,y)$ মোটেই~$y$-এর উপর নির্ভর করে না। $h$ যে আংশিক
গণনসাধ্য তা বোঝা কঠিন নয়: নিবেশ $x,y$-এ প্রথমে নিবেশ~$x$-এ
$\cfind{x}$ অপেক্ষকটিকে অনুকরণ করে~$h$ গণনা করি; এই গণনা থামলে
$h(x,y)$ মান~$0$ নির্গত করে থামে। তাই $h(x,y)$ শুধু
$\Zero(\umin{s}{T(x,x,s)})$, যেখানে $\Zero$ ধ্রুব শূন্য অপেক্ষক।

$s$-$m$-$n$ উপপাদ্য ব্যবহার করে এমন একটি আদিম পুনরাবৃত্ত অপেক্ষক~$k(x)$
পাওয়া যায়, যাতে প্রত্যেক $x$ ও~$y$-এর জন্য
\[
\cfind{k(x)}(y) \simeq
\begin{cases}
0 & \text{যদি $x \in K$} \\
\fundefined & \text{অন্যথায়}
\end{cases}
\]
% BN-SRC-155 TARGET-CORRECTION-BEGIN
\par\smallskip\noindent\textnormal{\small\textbf{উৎস-সংশোধন (BN-SRC-155):} হিমায়িত ইংরেজি প্রদর্শনে সম্ভাব্য অসংজ্ঞায়িত দুই পাশের মধ্যে সাধারণ “$=$” আছে। $x\notin K$ হলে উভয় পাশই অসংজ্ঞায়িত, তাই বাংলা পাঠে আংশিক অপেক্ষকের যুগপৎ সংজ্ঞায়িততার সমতা “$\simeq$” রাখা হয়েছে।} % BN-SRC-155 TARGET-CORRECTION-END
অতএব $x \in K$ হলে $\cfind{k(x)}$ সর্বত্র-সংজ্ঞায়িত, অন্যথায় সর্বত্র
অসংজ্ঞায়িত। তাই $k$, $K$ থেকে~$\fn{Tot}$-এ একটি হ্রাসকরণ।
\end{proof}

\begin{digress}
আসলে $\fn{Tot}$ এমনকি গণনসাধ্যভাবে তালিকায়নযোগ্যও নয়—“পাটীগাণিতিক
স্তরক্রমে” তার জটিলতা আরও উপরে। কিন্তু এই শক্তিশালী বক্তব্যটি এখানে
আমাদের বিবেচনার বিষয় নয়।
\end{digress}

\end{document}
